\chapter{Methodology}

\section{Development Approach}
BHUVAN was developed using an \textbf{Agile-Iterative} approach with sprint-based development. Given the solo-developer constraint, a modified Agile framework was adopted:

\begin{itemize}
    \item \textbf{Sprint duration:} 3--4 days per sprint.
    \item \textbf{Sprint planning:} Each sprint targeted one major subsystem (terrain, physics, UI, audio).
    \item \textbf{Continuous integration:} Vite's hot module replacement (HMR) provided instant feedback on code changes.
    \item \textbf{Iterative refinement:} Each subsystem was built as a minimum viable component, then iteratively improved (e.g., autopilot tuning required three iterations).
\end{itemize}

\section{Design Thinking Principles}
The development was guided by user-centric design thinking:

\begin{enumerate}
    \item \textbf{Empathize:} The primary users (interviewers, students, enthusiasts) want an instant, impressive experience---not a setup process.
    \item \textbf{Define:} The core problem is making terrain intelligence concepts tangible and interactive.
    \item \textbf{Ideate:} Multiple rendering approaches were evaluated (NeRF, 3DGS, procedural). Procedural generation was selected because it requires zero training data and runs on any hardware.
    \item \textbf{Prototype:} Each component was prototyped independently before integration.
    \item \textbf{Test:} Autopilot was tested across 50+ landing attempts; UI was tested at multiple resolutions.
\end{enumerate}

\section{Technology Stack Selection}

\begin{table}[h]
\centering
\caption{Technology Stack}
\label{tab:techstack}
\begin{tabular}{|p{3cm}|p{3.5cm}|p{6cm}|}
\hline
\textbf{Layer} & \textbf{Technology} & \textbf{Justification} \\
\hline
Build Tool & Vite 6.x & Sub-second HMR, ESM-native, minimal config \\
\hline
UI Framework & React 18 & Component-based architecture, hooks for state \\
\hline
3D Rendering & Three.js 0.169 & Industry-standard WebGL abstraction \\
\hline
3D React Binding & React Three Fiber 8.x & Declarative Three.js in JSX, useFrame hook \\
\hline
3D Utilities & Drei 9.x & OrbitControls, Stars, GLTF loaders \\
\hline
Post-Processing & @r3f/postprocessing & Bloom, vignette, chromatic aberration \\
\hline
Noise Generation & simplex-noise 4.x & Seeded 2D Simplex noise for terrain \\
\hline
Audio & Web Audio API & Procedural synthesis, zero file dependencies \\
\hline
Styling & Vanilla CSS & Maximum control, CSS custom properties \\
\hline
Fonts & Google Fonts & JetBrains Mono, Outfit \\
\hline
\end{tabular}
\end{table}

\section{Quality Assurance Strategy}
\begin{itemize}
    \item \textbf{Visual Testing:} Browser-based screenshot comparison after each sprint.
    \item \textbf{Physics Validation:} Autopilot tested against known analytical solutions for constant-gravity descent.
    \item \textbf{Performance Profiling:} Chrome DevTools Performance tab used to verify frame times $<$ 33ms.
    \item \textbf{Cross-Browser Testing:} Verified in Chrome 120, Firefox 121, and Edge 120.
\end{itemize}

\section{Tools and Collaboration}
\begin{itemize}
    \item \textbf{Version Control:} Git + GitHub for source code management.
    \item \textbf{IDE:} Visual Studio Code with ESLint and Prettier.
    \item \textbf{AI Assistance:} Claude, Gemini for code review and algorithm validation.
    \item \textbf{Documentation:} Overleaf for LaTeX report compilation.
\end{itemize}

\section{Project Timeline (Gantt Chart)}

\begin{figure}[h]
\centering
\begin{ganttchart}[
    hgrid, vgrid,
    x unit=0.7cm, y unit chart=0.6cm,
    bar/.style={fill=blue!40},
    milestone/.style={fill=red!60},
    title height=1,
    bar height=0.5
]{1}{15}
    \gantttitle{Days}{15} \\
    \gantttitlelist{1,...,15}{1} \\
    \ganttbar{Research \& Planning}{1}{2} \\
    \ganttbar{Terrain Engine}{3}{5} \\
    \ganttbar{Physics Engine}{4}{6} \\
    \ganttbar{3D Scene (R3F)}{5}{8} \\
    \ganttbar{Lander Model}{7}{8} \\
    \ganttbar{HUD \& UI}{8}{10} \\
    \ganttbar{Audio Engine}{9}{10} \\
    \ganttbar{Post-Processing}{10}{11} \\
    \ganttbar{Autopilot Tuning}{11}{12} \\
    \ganttbar{Camera Fix}{12}{13} \\
    \ganttbar{Testing \& Polish}{13}{14} \\
    \ganttbar{Documentation}{14}{15} \\
    \ganttmilestone{v1.0 Release}{15}
\end{ganttchart}
\caption{Project Development Timeline}
\label{fig:gantt}
\end{figure}
